\subsubsection{\Propagation}\label{ssec:designdims:propagation}

\begin{wrapfigure}{r}{0.37\linewidth}
\begin{lstlisting}[language=Py]
from trio import open_nursery

async def fail():
  raise Exception()

async def main():
  print("A")
  async with open_nursery() as n:
    n.start_soon(fail)
  print("B")
\end{lstlisting}
  \cprotect\caption{An example of \destruct \propagation. Python+Trio reraises exceptions when their associated nursery scope ends.}
  \label{fig:propagation:trio}
\end{wrapfigure}

One final point of divergence at a task's end of life is the handling of exceptions. In all surveyed languages with exceptions (Python, Swift, \Js, \CSharp), awaiting on an async function that has thrown an exception will re-raise the exception at the await point.\footnotemark{} However, what happens to an exception raised in a task that is never awaited? We call this the \propagation dimension.%
\cprotect\footnotetext{Instead of exceptions Rust uses the \code|Result| type to communicate errors. An \code|Err| variant is returned on await, which we consider a semantically equivalent design to re-raising an exception at the await point.}

Most languages \never propagate, meaning the exception generates no behavior outside its containing task, apart from perhaps printing to the console. Python+Asyncio and \Js will print a stack trace for an uncaught asynchronous exception. Node.js\ specifically will exit with a non-zero error code, although this latter behavior is not part of the ECMAScript specification so we still categorize it as a \never approach to \propagation.

Python+Trio is the only system which offers a different semantics, in part because tasks cannot be awaited, so the library must offer another mechanism for observing asynchronous exceptions. The example in \Cref{fig:propagation:trio} (shown in Python, not pseudocode) will print ``A'' but not ``B'' because the nursery will re-raise the exception originating in the \code|fail| function. We describe this as \destruct \propagation. A similar program in any other system will print ``A'' then ``B'', possibly with a stack trace printed in-between.

\paragraph{Design rationale.} The upside to the \destruct approach is that an application is less likely to ignore problematic exceptions. With the \never approach, if a developer does not carefully monitor their logs, a task can fail and cause confusing bugs where expected operations never occur. 

The downside is that the \destruct approach may run counter to either a developer's intuitions (i.e., exceptions aren't thrown away by the system) or their desired software architecture. Concurrency \abbr{api}s often seek to encapsulate errors at thread or task boundaries, so errors in one thread cannot crash the whole system. For example, in a backend web server with many concurrent connections, it is likely undesirable that one exceptional connection should crash the entire server. Either way, it seems sensible that that a straight-line asynchrony \abbr{api} should offer \emph{some} way of accessing information about unawaited exceptions, rather than only printing to standard error, or ignoring unawaited exceptions outright.
